% !TEX root = ../manuscript.tex

\section{Clustering~: identifier des profils de mode de vie}

\subsection{Objectif}

Le clustering est une méthode d'\textbf{apprentissage non supervisé}~: il n'y a pas de variable cible. L'objectif est de regrouper les individus selon leur similarité pour identifier des \textbf{profils types de mode de vie}~\cite{macqueen1967}. La variable \texttt{disease\_risk} est exclue du clustering.

\subsection{Préparation des données}

Les 15 variables explicatives sont standardisées (z-score) avant d'appliquer les algorithmes. K-Means et CAH reposent sur la \textbf{distance euclidienne}~\cite{macqueen1967}~: sans standardisation, \texttt{daily\_steps} (valeurs jusqu'à 20~000) dominerait mécaniquement \texttt{sleep\_hours} (3 à 10), et le clustering serait réduit à une partition par cette seule variable. Le z-score garantit une contribution équitable de chacune des 15 dimensions~\cite{hastie2009}.

La CAH étant de complexité $\mathcal{O}(n^2)$ (voir note sur le SVM au chapitre~3), elle est appliquée sur un sous-échantillon aléatoire de \textbf{5~000 individus} pour rester dans des temps de calcul raisonnables.

\subsection{Sélection du nombre de clusters}

\textbf{Méthode du coude}~: on trace l'\textbf{inertie intra-cluster} $W(k)$ en fonction de $k$. L'inertie mesure la dispersion totale à l'intérieur des clusters~:
\begin{equation}
W(k) = \sum_{j=1}^{k}\sum_{\mathbf{x} \in C_j}\|\mathbf{x} - \boldsymbol{\mu}_j\|^2
\end{equation}
où $C_j$ désigne le $j$-ème cluster et $\boldsymbol{\mu}_j$ son centroïde (vecteur des moyennes). Plus $k$ augmente, plus l'inertie diminue --- mais à partir d'un certain $k$, le gain devient marginal~: c'est le "coude" de la courbe, au-delà duquel ajouter des clusters ne compense plus la complexité.

\textbf{Coefficient de silhouette}~\cite{rousseeuw1987}~: défini au chapitre~1, il mesure la qualité de séparation des clusters. Pour rappel~:
\begin{equation}
s(i) = \frac{b(i) - a(i)}{\max(a(i),\, b(i))}
\end{equation}
où $a(i)$ est la distance moyenne de l'individu $i$ aux autres membres de son cluster, et $b(i)$ la distance moyenne au cluster voisin le plus proche. Une valeur proche de 1 indique une bonne séparation.

\begin{table}[h]
\centering
\caption{Inertie et silhouette selon $k$ (K-Means)}
\label{tab:kmeans_k}
\begin{tabular}{rrr}
\toprule
$k$ & Inertie & Silhouette \\
\midrule
2 & 1~346~768 & 0.0979 \\
\textbf{3} & \textbf{1~238~076} & \textbf{0.1030} \\
4 & 1~188~303 & 0.0967 \\
5 & 1~152~260 & 0.0886 \\
10 & 1~031~548 & 0.0744 \\
\bottomrule
\end{tabular}
\end{table}

$k=3$ maximise la silhouette et correspond au coude de la courbe d'inertie. \textbf{$k=3$ est retenu.}

\begin{figure}[h]
\centering
\includegraphics[width=0.85\textwidth]{kmeans_selection_k.png}
\caption{Inertie intra-cluster et coefficient de silhouette selon $k$ (K-Means)}
\label{fig:kmeans_k}
\end{figure}

\subsection{K-Means}

L'algorithme K-Means~\cite{macqueen1967} alterne affectation au centroïde le plus proche et mise à jour des centroïdes~: $\boldsymbol{\mu}_j = \frac{1}{|C_j|}\sum_{\mathbf{x} \in C_j} \mathbf{x}$, jusqu'à convergence.

\begin{table}[h]
\centering
\caption{Profils des clusters K-Means ($k=3$)}
\label{tab:kmeans_profiles}
\begin{tabular}{llrrl}
\toprule
Cluster & Taille & IMC moyen & Syst. (mmHg) & Profil \\
\midrule
0 & 32~754 & 35.0 (obèse)    & 142 (hypertendu)  & Obèse hypertendu \\
1 & 27~673 & 29.1 (surpoids) & 115 (normotendu)  & Surpoids normotendu \\
2 & 39~573 & 24.0 (normal)   & 142 (hypertendu)  & Poids normal hypertendu \\
\bottomrule
\end{tabular}
\end{table}

\begin{quote}
\textbf{Interprétation~:} K-Means partitionne la population selon l'\textbf{IMC} et la \textbf{pression artérielle systolique}. Ce résultat s'explique par la sensibilité de K-Means aux variables à grande variance~: dans notre dataset, \texttt{bmi} (écart-type = 6.3) et \texttt{systolic\_bp} (écart-type = 26.0) ont les dispersions les plus larges après standardisation, ce qui les rend naturellement "dominantes" pour la distance euclidienne. Le cluster~1 regroupe les individus les plus sains (surpoids mais normotendus), tandis que les clusters~0 et~2 correspondent à deux profils à risque cardiovasculaire distincts selon la classification OMS~\cite{who2000} et ESC~\cite{esc2018}.
\end{quote}

\subsection{Classification Hiérarchique Ascendante (CAH)}

La CAH construit un \textbf{dendrogramme} en fusionnant progressivement les clusters les plus proches. Au départ, chaque individu est son propre cluster~; à chaque étape, on fusionne les deux clusters dont la fusion est la "moins coûteuse"~\cite{hastie2009}.

La \textbf{méthode de liaison Ward} définit le coût de fusion de $A$ et $B$ comme l'augmentation de variance intra-cluster~:
\begin{equation}
\Delta(A, B) = \frac{|A|\cdot|B|}{|A|+|B|}\|\boldsymbol{\mu}_A - \boldsymbol{\mu}_B\|^2
\end{equation}
Ward minimise ce coût à chaque étape, favorisant des clusters \textbf{compacts et équilibrés}~--- d'où sa sensibilité aux sous-groupes comportementaux homogènes dans nos données.

\begin{table}[h]
\centering
\caption{Profils des clusters CAH ($k=3$, $n=5~000$)}
\label{tab:cah_profiles}
\begin{tabular}{llrrl}
\toprule
Cluster & Taille & Fumeurs & Syst. (mmHg) & Profil \\
\midrule
0 & 2~928 & 2\%           & 142 & Non-fumeurs hyperten dus \\
1 & 1~407 & 18\%          & 115 & Fumeurs modérés, normotendus \\
2 &   665 & \textbf{100\%}& 142 & Fumeurs intenses, hyperten dus \\
\bottomrule
\end{tabular}
\end{table}

\begin{quote}
\textbf{Interprétation~:} La CAH isole un \textbf{cluster de fumeurs intenses} (100\% de fumeurs + hypertendus). Ce résultat s'explique par la nature de l'algorithme~: Ward cherche des sous-groupes \textit{homogènes}, y compris sur des variables binaires comme \texttt{smoker}. K-Means, guidé par les grandes variances numériques, n'aurait pas isolé ce groupe de 665 individus --- il les aurait dilués dans les clusters dominés par l'IMC et la pression. Les deux méthodes sont donc \textbf{complémentaires}~: K-Means donne une vue globale par dimensions biologiques continues, la CAH révèle des sous-groupes comportementaux fins. Le profil tabagisme + hypertension est un double facteur de risque cardiovasculaire bien documenté~\cite{booth2012}.
\end{quote}

\subsection{Évaluation et visualisation}

\textbf{Visualisation par ACP (Analyse en Composantes Principales).} Nos données vivent dans un espace à 15 dimensions~: il est impossible de les représenter directement. L'ACP projette ces 15 dimensions sur les 2 axes qui capturent le maximum de variance, permettant une visualisation 2D des clusters. Les deux premières composantes expliquent \textbf{24.3\%} de la variance totale~: la projection 2D "voit" seulement 24.3\% de l'information originale. Si les clusters semblent peu séparés sur le graphique, ce n'est pas la preuve qu'ils n'existent pas --- ils se distinguent dans les 15 dimensions d'origine, pas dans la réduction 2D. Avec 15 variables quasi-indépendantes (confirmé par l'EDA), aucune direction ne concentre beaucoup de variance~: 24.3\% en 2D est donc attendu.

\textbf{Pourquoi les clusters ont-ils une forme rectangulaire dans la projection ACP~?} Sur le graphique ACP (figure~\ref{fig:kmeans_pca}), les clusters apparaissent comme des \textbf{blocs rectangulaires} plutôt que des nuages sphériques --- résultat directement lié à la nature des variables du dataset.

La première composante principale (PC1, axe horizontal) est fortement influencée par \texttt{bmi\_cat}, une variable ordinale à \textbf{4 valeurs discrètes} (0 = sous-poids, 1 = normal, 2 = surpoids, 3 = obèse). Cette discrétisation crée des \textbf{colonnes verticales}~: les individus s'alignent en 4 bandes distinctes plutôt que de se répartir continûment sur l'axe. La deuxième composante (PC2, axe vertical) est dominée par \texttt{hypertension}, une variable \textbf{binaire} (0 ou 1)~: une variable binaire ne peut prendre que deux valeurs, ce qui produit un \textbf{saut net haut/bas} dans la projection, divisant les individus entre normotendus (bas) et hypertendus (haut).

Par ailleurs, les distributions quasi-uniformes de toutes les variables continues (chapitre~2) créent des nuages de points dispersés sans gradient de densité visible~: les points se répartissent uniformément au lieu de former des ellipses denses au centre. Le résultat final est une grille de rectangles dont la forme reflète la structure discrète du dataset plutôt qu'une vraie géométrie en clusters.

\textbf{Indice de Davies-Bouldin~\cite{davies1979}.} Cet indice mesure la qualité du clustering en calculant, pour chaque cluster $i$, le ratio entre la dispersion interne et la distance entre centroïdes~:
\begin{equation}
DB = \frac{1}{k}\sum_{i=1}^{k}\max_{j \neq i}\left(\frac{\sigma_i + \sigma_j}{d(\boldsymbol{\mu}_i, \boldsymbol{\mu}_j)}\right)
\end{equation}
où $\sigma_i$ est la dispersion moyenne des points du cluster $i$ autour de son centroïde, et $d(\boldsymbol{\mu}_i, \boldsymbol{\mu}_j)$ la distance entre les centroïdes $i$ et $j$. \textbf{Un DB proche de 0 est idéal}~: les clusters sont compacts (faible $\sigma$) et bien séparés (grande distance). Un DB élevé indique des clusters éparpillés ou trop proches les uns des autres. Contrairement à la silhouette qui évalue chaque point individuellement, Davies-Bouldin évalue chaque paire de clusters~: les deux métriques sont complémentaires.

\begin{table}[h]
\centering
\caption{Comparaison K-Means vs CAH}
\label{tab:clustering_eval}
\begin{tabular}{lrrl}
\toprule
Métrique & K-Means & CAH & Interprétation \\
\midrule
Silhouette      & 0.103 & 0.105 & Faible~--- clusters peu séparés \\
Davies-Bouldin  & 2.655 & 2.721 & Élevé~--- chevauchement entre clusters \\
Variance ACP (2D) & \multicolumn{2}{c}{24.3\%} & Projection insuffisante (15 dims $\to$ 2D) \\
\bottomrule
\end{tabular}
\end{table}

\begin{figure}[h]
\centering
\includegraphics[width=0.85\textwidth]{kmeans_pca.png}
\caption{Visualisation des clusters K-Means ($k=3$) par projection ACP (2 composantes, 24.3\% de variance)}
\label{fig:kmeans_pca}
\end{figure}

Malgré des silhouettes faibles ($\approx 0.10$), les deux méthodes révèlent des structures interprétables médicalement~\cite{booth2012, who2000}~: obèse hypertendu (K-Means), fumeurs intenses + hypertendus (CAH). La faible silhouette reflète l'absence de groupes naturels francs dans des données synthétiques générées sans structure causale --- pas une erreur méthodologique.

% ─────────────────────────────────────────────────────────────
\section{Interprétabilité des modèles}

\subsection{Objectif}

Cette section pose une question précise~: \textbf{les outils d'interprétabilité convergent-ils vers les mêmes variables influentes, même en l'absence de signal~?} La réponse est elle-même un résultat analytique.

\subsection{Outils utilisés}

\textbf{Coefficients des modèles linéaires}~: un coefficient $w_j > 0$ indique qu'une augmentation de la variable $j$ est associée à une augmentation de la prédiction~\cite{james2021}. Leur magnitude reflète l'importance relative de chaque variable --- sous réserve que les variables soient standardisées.

\textbf{Feature Importance (impureté de Gini)}~\cite{breiman2001}~: pour les forêts aléatoires, l'importance de la variable $j$ agrège la réduction d'impureté de Gini qu'elle apporte à travers tous les nœuds de tous les arbres. Plus une variable est utilisée haut dans les arbres pour des décisions qui séparent bien les classes, plus son importance est élevée. La limite connue de cet indicateur~: il favorise les variables à haute cardinalité ou grande variance --- ce qui, dans notre cas, explique que \texttt{systolic\_bp} domine systématiquement.

\textbf{SHAP}~: défini au chapitre~1 (formule des valeurs de Shapley) et utilisé ici comme complément visuel, il attribue à chaque variable une contribution signée à chaque prédiction individuelle. Les formules ne sont pas redéfinies ici.

\subsection{Résultats, convergence et analyse critique}

Les cinq méthodes d'interprétabilité (coefficients LR, Ridge, feature importance RF, XGBoost, SHAP) sont appliquées et leurs top~5 variables comparés.

\begin{quote}
\textbf{Analyse critique~:} Les cinq méthodes ne s'accordent pas sur un classement stable~--- \textbf{cette divergence est elle-même un résultat}~\cite{rudin2019}~: sans signal prédictif réel, chaque outil mesure du bruit aléatoire et s'accroche à des corrélations parasites différentes. La seule variable présente dans les cinq méthodes est \texttt{systolic\_bp}~: dans la littérature épidémiologique réelle, la pression artérielle est effectivement l'un des facteurs de risque cardiovasculaire les mieux documentés~\cite{esc2018}~--- mais dans nos données, cette convergence résulte de sa grande variance, pas d'une vraie relation causale.

Sur des données réelles~\cite{booth2012, who2000}, on s'attendrait à voir la pression artérielle, le cholestérol LDL, le tabagisme, l'IMC et l'activité physique dominer clairement les classements, avec des coefficients significativement différents de zéro.
\end{quote}

\begin{figure}[h]
\centering
\includegraphics[width=0.75\textwidth]{lr_coefficients.png}
\caption{Coefficients de la régression logistique (variables standardisées)}
\label{fig:lr_coef}
\end{figure}

\begin{figure}[h]
\centering
\includegraphics[width=0.75\textwidth]{rf_feature_importance.png}
\caption{Feature importances (impureté de Gini) du Random Forest}
\label{fig:rf_fi}
\end{figure}

\begin{figure}[h]
\centering
\includegraphics[width=0.85\textwidth]{shap_summary.png}
\caption{SHAP summary plot (XGBoost) --- complément visuel d'interprétabilité}
\label{fig:shap}
\end{figure}

\begin{quote}
\textbf{Comment lire le SHAP summary plot~?} Chaque \textbf{ligne} correspond à une variable, triée de haut en bas par importance décroissante (impact moyen en valeur absolue). Sur chaque ligne, chaque \textbf{point} représente un individu. L'\textbf{axe horizontal} indique l'impact SHAP sur la prédiction~: un point à droite ($x > 0$) signifie que la variable pousse la prédiction vers 1 (\texttt{disease\_risk} élevé)~; un point à gauche ($x < 0$) la pousse vers 0~; en $x = 0$, aucun impact. La \textbf{couleur} encode la valeur réelle de la variable~: rouge = valeur élevée, bleu = valeur faible.

Sur des données avec un signal réel, on verrait un gradient net~: par exemple, si une pression artérielle élevée (rouge) est un facteur de risque, les points rouges se concentreraient à droite et les bleus à gauche. Ici, les couleurs sont \textbf{mélangées aléatoirement} autour de zéro~: des points rouges et bleus occupent les mêmes positions horizontales, sans aucun pattern. Cela confirme qu'aucune variable n'a d'impact systématique~--- SHAP détecte du bruit plutôt qu'un signal.
\end{quote}

Le tableau suivant synthétise la présence de chaque variable dans le top~5 des cinq méthodes~:

\begin{table}[h]
\centering
\caption{Variables dans le top 5 de chaque méthode d'interprétabilité}
\label{tab:interpretability}
\begin{tabular}{lcccccc}
\toprule
Variable & LR & RF & XGB & SHAP & Ridge & Score \\
\midrule
\texttt{systolic\_bp}   & \checkmark & \checkmark & \checkmark & \checkmark & \checkmark & \textbf{5/5} \\
\texttt{daily\_steps}   & \checkmark & \checkmark & ---        & \checkmark & ---        & 3/5 \\
\texttt{family\_history}& \checkmark & ---        & \checkmark & ---        & \checkmark & 3/5 \\
\texttt{bmi}            & ---        & \checkmark & ---        & \checkmark & ---        & 2/5 \\
\bottomrule
\end{tabular}
\end{table}
